% Nouveau EGA V (cote 25): archive cote 25, site page 82 (PDF page 83).
% typescript. Draft compared: none.
% AI-transcribed by Claude Opus 5.5 (claude-opus-5-5, adaptive thinking, effort high), Message Batches API, batch ega25b (msgbatch_01A7n5YuXubUPNK1zsepi8RG), 2026-11-13; checked against both supplied scan parts; see CHECK_DECISIONS.tsv.
\scanpage{25/82}
\mspage{\hand{77}}

de base général. Signalons la suivante (sans démonstration): Soit $f:X\hand{\to} S$
un morphisme projectif et plat, \var{\xout{\unsure{sur S tel que X}}}{}{typewriter strike-out}{ega25_25_082-v0001} à fibres géométriques
\var{\xout{géomé}}{}{typewriter strike-out}{ega25_25_082-v0002} irréductibles et $(R_2)$ , $S'$ un sous-\var{\inserted[typed]{pré}}{pré}{insertion}{ega25_25_082-v0003}schéma de $X$ fini sur $S$,
$s\hand{\in} S$, supposons que pour tout $x\hand{\in} S'$ au dessus de $s$, $X_s$ soit lisse sur
$k(s)$ en $x$ . Alors il existe \unsure{un} voisinage ouvert $U$ de $s$, et un voisinage
ouvert $V$ de $S'|U$ dans $X|U$ , tels que $V\hand{\longrightarrow} U$ soit poly-élémentaire.
Si $Y$ est un sous-préschéma fermé de $X$\var{\inserted[typed]{ne rencontrant pas $S'$ et}}{ ne rencontrant pas $S'$ et}{insertion}{ega25_25_082-v0004} tel que l'ensemble $Z$ des points
en lesquels $Y$ n'est pas lisse sur $S$ \var{\xout{ait des fibres de}}{}{typewriter strike-out}{ega25_25_082-v0005} \var{\deleted{soit tel que}}{}{deletion}{ega25_25_082-v0006} \var{\inserted{satisfasse}}{satisfasse}{insertion}{ega25_25_082-v0007}
dim $Z_s$ \hand{$\leq$} dim $X_s$ - 2 , alors on peut ci-dessus prendre $V$ contenu dans
$X$ - $Y$.

c) (A être éventuellement inclus dans l'énoncé de 13.2.). Soit
avec les notations de 13.2. $\hand{\unsure{\phi|\emptyset}}$ une partie finie de $X$ formée de points
lisses de $X$, et supposons que $\hand{\unsure{\phi|\emptyset}}$ soit contenu dans un \var{\xout{\illegible{1}}}{}{typewriter strike-out}{ega25_25_082-v0008} ouvert affine
de $X$. Alors $\hand{\unsure{\phi|\emptyset}}$ admet un systême fondamental de voisinages ouverts
polyélémentaires. \var{\inserted[typed]{Evidemment on peut supposer $\hand{\unsure{\phi|\emptyset}}$ formé de points \uline{fermés}.}}{Evidemment on peut supposer $\hand{\unsure{\phi|\emptyset}}$ formé de points \uline{fermés}.}{insertion}{ega25_25_082-v0009} La démonstration est essentiellement la même, sauf
qu'on formule 13.3. sous une forme légèrement différente, savoir: il
existe une sous-variété linéaire $C$ de $P^r$ de codimension $n$, ne rencon\unsure{t|-}%
trant pas $\hand{\unsure{\phi|\emptyset}}\hand{\cup} Y$ , et telle que pour tout $x_i\hand{\in}\hand{\unsure{\phi|\emptyset}}$\var{\inserted[typed]{$(\overset{\hand{-}}{k})$}}{$(\overset{\hand{-}}{k})$}{insertion}{ega25_25_082-v0010}, la sous-variété
linéaire $L_i$ de codimension $n$-1 engendrée par $C$ et $x_i$ ait les propriétés
a) , b) , c) de 13.3. Pour vérifier ce point, on note qu'il suffit de
vérifier que la $C$ \uline{générique} a les propriétés envisagées, or \var{\inserted{pour une telle $C$}}{pour une telle $C$}{insertion}{ega25_25_082-v0011} \var{\deleted{alors}}{}{deletion}{ega25_25_082-v0012}
\var{\xout{\unsure{les $L$}}}{}{typewriter strike-out}{ega25_25_082-v0013}\unsure{\hand{,}|$_i$} chaque $L_i$ est génér\unsure{ique} parmi les $L$ de codimension $n$ passant
par $x_i$ , de sorte qu'on peut appliquer 13.3. sous la forme initiale
( ou du moins sous la forme où nous l'avions prouvée, qui était: toute
$L_o$ \uline{assez générale} passant par $x$ a les propriétés a),b),c) ).

e) Une des raisons de l'intérêt de 13.2. est la structure``topolo-%
gique'' particulièrement simple des \var{\xout{\unsure{ouverts}}}{}{typewriter strike-out}{ega25_25_082-v0014} \var{\inserted[typed]{schémas alg.}}{schémas alg.}{insertion}{ega25_25_082-v0015} élémentaires\var{\inserted{$U$}}{ $U$}{insertion}{ega25_25_082-v0016}. Par exemple,

\mspage{\hand{82}}
